\documentclass[ijoc,nonblindrev]{informs4} 

% Setup for Electronic Companion
\usepackage{xr}
\externaldocument{BESO_JoC} % Link to your main paper's aux file

% Prefix numbering
\renewcommand{\thesection}{EC.\arabic{section}}
\renewcommand{\theequation}{EC.\arabic{equation}}
\renewcommand{\thefigure}{EC.\arabic{figure}}
\renewcommand{\thetable}{EC.\arabic{table}}

\begin{document}
	
%	\TITLE{e-companion to: Bootstrap Enhanced Scenario Optimization, a Case Study in Two-Echelon Logistics}
%	\AUTHOR{Livio Fenga, Vittorio Maniezzo}
%	
%	\maketitle
	
	\section*{Electronic Companion}
	\ECHead{Electronic Companion to "Bootstrap Enhanced Scenario Optimization, a Case Study in Two-Echelon Logistics"}
	
	% Your supplemental content starts here
% Appendix here
% Options are (1) APPENDIX (with or without general title) or
%             (2) APPENDICES (if it has more than one unrelated sections)
% Outcomment the appropriate case if necessary
%
% \begin{APPENDIX}{<Title of the Appendix>}
	% \end{APPENDIX}
%
%   or
%
\begin{APPENDICES}
	\section{EC-Bootstrap sets error costs} \label{app:errcost}
	
	This section reports the absolute values of the various cost functions, which were calculated using the forecast errors obtained by averaging the corresponding errors for each of the 52 test case series described in the paper.
	
	\begin{table}[h]
		\centering
		\renewcommand{\arraystretch}{0.8}
		\begin{tabular}{ccccccc}
			model & bias & MAPE & ME & MAE & MPE & RMSE \\
			\hline
			fcast\_50   &-0.27 & 0.09 &-1.02 & 1.63 &-0.05 &  2.12 \\
			fcast\_avg  & 0.02 & 0.07 & 0.07 & 1.25 & 0.01 &  1.68 \\
			yar         & 0.36 & 0.34 & 1.36 & 5.86 & 0.07 &  9.75 \\
			yhw         & 0.63 & 0.34 & 2.41 & 5.99 & 0.13 & 10.59 \\
			ysvm        & 0.69 & 0.36 & 2.66 & 6.20 & 0.14 & 10.99 \\
			ylstm       & 0.56 & 0.35 & 2.13 & 6.01 & 0.11 & 10.41 \\
			ymlp        & 0.47 & 0.36 & 1.79 & 6.26 & 0.09 & 10.60 \\
			yrf         & 1.50 & 0.44 & 5.73 & 7.79 & 0.31 & 13.54 \\
			yxgb        & 2.28 & 0.58 & 8.74 &10.12 & 0.49 & 17.05 \\
			yarima      & 0.44 & 0.34 & 1.67 & 5.80 & 0.09 &  9.77 \\
			\hline
		\end{tabular}
		\caption{Error costs for the 75 series boostset.}
		\label{table:app75}
	\end{table}
	
	\begin{table}[h]
		\centering
		\renewcommand{\arraystretch}{0.8}
		\begin{tabular}{ccccccc}
			model & bias & MAPE & ME & MAE & MPE & RMSE \\
			\hline
			fcast\_50   & -0.56 & 0.09 & -1.08 & 1.59 & -0.06 & 2.25  \\
			fcast\_avg  &  0.02 & 0.07 &  0.05 & 1.18 &  0.01 & 1.65  \\
			yar         & -0.59 & 0.25 & -1.13 & 4.36 & -0.07 & 6.31  \\
			yhw         & -0.13 & 0.24 & -0.26 & 4.33 & -0.02 & 6.39  \\
			ysvm        & -0.06 & 0.24 & -0.11 & 4.18 & -0.01 & 6.39  \\
			ylstm       & -0.30 & 0.24 & -0.58 & 4.15 & -0.04 & 6.34  \\
			ymlp        & -0.47 & 0.24 & -0.90 & 4.22 & -0.05 & 5.99  \\
			yrf         &  1.27 & 0.28 &  2.44 & 5.01 &  0.14 & 7.48  \\
			yxgb        &  2.52 & 0.38 &  4.85 & 6.52 &  0.28 & 9.46  \\
			yarima      & -0.40 & 0.25 & -0.76 & 4.44 & -0.05 & 6.83  \\
			\hline
		\end{tabular}
		\caption{Error costs for the 125 series boostset.}
		\label{table:app125}
	\end{table}
	
	\begin{table}[h]
		\centering
		\renewcommand{\arraystretch}{0.8}
		\begin{tabular}{ccccccc}
			model & bias & MAPE & ME & MAE & MPE & RMSE \\
			\hline
			fcast\_50  & -0.04 & 0.07 & -0.07 & 1.22 & 0     & 1.51 \\
			fcast\_avg & 0.99  & 0.12 & 1.90  & 2.04 & 0.12  & 2.76 \\
			yar        & -0.22 & 0.35 & -0.42 & 5.84 & -0.03 & 7.64 \\
			yhw        & 0.25  & 0.34 & 0.48  & 5.75 & 0.02  & 7.90 \\
			ysvm       & 0.30  & 0.33 & 0.58  & 5.66 & 0.03  & 7.76 \\
			ylstm      & 0.10  & 0.34 & 0.20  & 5.76 & 0     & 7.74 \\
			ymlp       & -0.19 & 0.34 & -0.36 & 5.69 & -0.02 & 7.57 \\
			yrf        & 1.72  & 0.37 & 3.31  & 6.39 & 0.18  & 9.66 \\
			yxgb       & 3.01  & 0.46 & 5.78  & 7.96 & 0.33  & 11.83\\
			yarima     & -0.04 & 0.34 & -0.08 & 5.69 & -0.01 & 7.63 \\
			\hline
		\end{tabular}
		\caption{Error costs for the 175 series boostset.}
		\label{table:app175}
	\end{table}
	
	The functions reported in the columns are: bias, Means Absolute Percentage Error (MAPE), Mean Error (ME), Mean Absolute Error (MAE), Mean Percentage Error (MPE) and Root Mean Square Error (RMSE). Table \ref{table:app75} gives the cost values for the case of a boost set composed of 75 series, table \ref{table:app125} for the 125 series and table \ref{table:app175} for the 175 series sets.
	
	All tables show significantly consistent results across all boost set sizes.
	
	\section{EC-Extended deterministic benchmark set} \label{app:artificial}
	
	All analyses reported in sections \ref{Subsec:predictive} and \ref{Subsec:prescriptive} refer to the case study at the heart of this paper, but the interest of the method is not limited to this particular case study as it can be effective on any instance of this 2-echelon supply chain problem. To test the generality of the results, we generated a benchmark set with varying dimensions and coefficients. The study was conducted only with respect to the deterministic setting, and is therefore significant for assessing the computational complexity of the split-allocation problem as a deterministic combinatorial optimization problem.
	
	We generated 5 instances for each configuration of number of clients ($n$), number of servers ($m$), and number of splittable clients. The travel costs were obtained by randomly selecting $m$ rows and $n$ columns from a large distance matrix, whose values were computed as the actual road distances between points located in the same area as the test case. Server inventory costs were inversely proportional to the average travel cost from the server to the clients and client requests were inversely proportional to the average travel cost from the client to the servers. Note that the distance matrix is asymmetric, as it was computed within a city area. The number $n_b$ of splittable clients was progressively increased as the number $n_q$ of servers with nonzero inventory costs. Specifically, the range of the defining parameters were: $n$: \{50, 52, 100, 200, 300\}, $m$: \{4, 5, 10, 25, 50\}, $n_b$: \{0, 4, 5, 10, 25, 50\}, $n_q$: \{0, 1, 2, 3, 4, 5\}.
	
	Not all combinations of values were generated, the larger values were used only for instances with higher number of clients. In particular, the values 4 and 52 were only used to generate instances with the same dimensionality as the case study, in order to check the generality of the results obtained. The first instance of this set was kept to be the case study one.
	
	The computational results in the case of no costly server, reported as averages over the 5 generated instances, are summarized in table \ref{table:extset}, whose columns show:
	\begin{itemize}
		\item $n$, $m$: number of clients and number of servers;
		\item $n_b$: maximum number of splittable clients;
		\item $ncols$, $nrows$: number of columns and number of rows in formulation FD;
		\item $gaplb$: average percentage gap between the linear relaxation lower bound and the cost of the best solution found;
		\item $gapfinal$: average percentage gap between the best lower bound at the end of the search (smallest cost of an unexpanded node) and the cost of the best solution found;
		\item $ninf$: number of instances for which no feasible solution was found;
		\item $nopt$: number of instances solved to proven optimality;
		\item $tcpu$: average CPU time for solving the instance, time to optimality or 3600 if optimality could not be proven;
	\end{itemize}
	
	The table shows two trends. One, with respect to instance size, shows an expected increase in complexity. While instances the size of our test case are easily solved to optimality by modern solvers, increasing the size, as defined by the number of variables and the number of constraints, quickly leads to instances that could not be solved to optimality within the 3600-second time limit. 
	The number of splittable clients has an important effect on the complexity of the instance. Instances of the same absolute size, but with more splittable clients, usually have a looser linear relaxation (LP) bound and have more difficulty reaching optimality, as can be seen by the decreasing number of instances out of the 5 from each group that could be solved to optimality.
	
	\begin{table}
		\centering
		\renewcommand{\arraystretch}{0.8}
		\begin{tabular}{cccccccccc}
			$n$ & $m$ & $n_b$ & $ncols$ & $nrows$ & $gaplb$ & $gapfinal$ & $ninf$ & $nopt$ & $tcpu$\\
			\hline
			52 &  4 &  0 &   416 &   316 & 0.21 & 0.01 & 0 & 5 &    0  \\
			\hline
			50 &  5 &  0 &   500 &   355 & 0.39 & 0.00 & 0 & 5 &    3  \\
			50 &  5 &  5 &   500 &   355 & 0.40 & 0.00 & 0 & 5 &    4  \\
			50 & 10 &  0 &  1000 &   610 & 1.32 & 0.00 & 0 & 5 &  383  \\
			50 & 10 &  5 &  1000 &   610 & 1.34 & 0.00 & 0 & 5 &  378  \\
			50 & 10 & 10 &  1000 &   610 & 1.35 & 0.13 & 0 & 4 & 1150  \\
			\hline
			100 & 10 &  0 &  1000 &  1210 & 0.50 & 0.12 & 0 & 2 & 1920  \\
			100 & 10 &  5 &  1000 &  1210 & 0.54 & 0.05 & 0 & 3 & 1514  \\
			100 & 10 & 10 &  1000 &  1210 & 0.57 & 0.09 & 0 & 3 & 2063  \\
			100 & 50 &  0 & 10000 &  5250 & 5.23 & 3.06 & 0 & 0 & 3600  \\
			100 & 50 &  5 & 10000 &  5250 & 4.77 & 2.70 & 0 & 0 & 3600  \\
			100 & 50 & 10 & 10000 &  5250 & 4.55 & 2.42 & 0 & 0 & 3600  \\
			100 & 50 & 25 & 10000 &  5250 & 4.52 & 2.61 & 0 & 0 & 3600  \\
			\hline
			200 & 10 &  0 &  4000 &  2410 & 0.17 & 0.05 & 0 & 2 & 2369  \\
			200 & 10 &  5 &  4000 &  2410 & 0.17 & 0.09 & 0 & 1 & 2893  \\
			200 & 10 & 10 &  4000 &  2410 & 0.18 & 1.01 & 0 & 0 & 3600  \\
			200 & 50 &  0 & 20000 & 10450 & 1.66 & 1.26 & 0 & 0 & 3600  \\
			200 & 50 & 10 & 20000 & 10450 & 1.68 & 1.29 & 0 & 0 & 3600  \\
			200 & 50 & 25 & 20000 & 10450 & 1.71 & 1.37 & 0 & 0 & 3600  \\
			\hline
			300 & 10 &  0 &  6000 &  3610 & 0.09 & 0.03 & 0 & 3 & 1798 \\
			300 & 10 &  5 &  6000 &  3610 & 0.09 & 0.07 & 0 & 0 & 3600 \\
			300 & 10 & 10 &  6000 &  3610 & 1.02 & 0.47 & 0 & 0 & 3600 \\
			300 & 50 &  0 & 30000 & 15650 & 0.94 & 0.73 & 0 & 0 & 3600  \\
			300 & 50 & 10 & 30000 & 15650 & 0.93 & 0.76 & 0 & 0 & 3600  \\
			300 & 50 & 25 & 30000 & 15650 & 0.96 & 0.84 & 0 & 0 & 3600  \\
			\hline
		\end{tabular}
		\caption{Extended benchmark set, deterministic case, no costly server}
		\label{table:extset}
	\end{table}
	
	Note that the optimality gap of the LP bound is usually small, except for the $n$=100 and $m$=50 instances, but the search still struggles to close this gap. Interestingly, the LP bound gets tighter as $n$ increases, but this does not make the search any easier.
	The least-cost solution when optimality was not proven was usually found within the first 600 seconds of CPU time. However, we note that designing efficient heuristics for this problem is beyond the scope of this research. We could have tweaked the solver configuration to favor heuristic search at the expense of efficiency in proving optimality, or designed custom heuristics, but the focus here is on assessing the fitness landscape when road network costs are involved.
	
	Table \ref{table:extsetq} reports results about tests varying the number of costly servers, the columns show:
	\begin{itemize}
		\item $n$, $m$: number of clients and number of servers;
		\item $ncols$, $nrows$: number of columns and number of rows in formulation FD;
		\item lb0, t0: linear relaxation bound and CPU time in the case of no costly servers;
		\item lb1, t1: linear relaxation bound and CPU time in the case of 1 costly server;
		\item lb2, t2: linear relaxation bound and CPU time in the case of 2 costly servers;
		\item lb3, t3: linear relaxation bound and CPU time in the case of 3 costly servers;
		\item lb4, t4: linear relaxation bound and CPU time in the case of 4 costly servers;
		\item lb5, t5: linear relaxation bound and CPU time in the case of 5 costly servers;
	\end{itemize}
	
	The table reports results on two sets of instances. The top 4 rows contain aggregate results on a set of instances, 5 for each configuration, where inventory costs are directly proportional to average distances, while the bottom 4 rows have costs inversely proportional to average distances.
	
	\begin{table}
		\addtolength{\tabcolsep}{-3pt}
		\footnotesize 
		\centering
		\renewcommand{\arraystretch}{1.2}
		\begin{tabular}{cccc|cc|cc|cc|cc|cc|cc}
			n & m & ncols & nrows & lb0 & t0 & lb1 & t1 & lb2 & t2 & lb3 & t3 & lb4 & t4 & lb5 & t5 \\
			\hline
			52 & 4  & 416 & 316 & 0.86 & 0.2 & 0.17 & 0.81 & 0.08 & 0.22 & 0.01 & 0.02 & 0.00 & 0.02 & - & - \\
			50 & 5  & 500 & 355 & 1.33 & 0.06 & 2.77 & 0.06 & 1.46 & 0.13 & 0.11 & 14.31 & 0.06 & 7 & 0.05 & 0.93 \\
			100 & 5 & 1000 & 705 & 0.41 & 0.24 & 0.09 & 6.39 & 0.06 & 3.2 & 0.01 & 0.06 & 0.02 & 4.8 & 0.02 & 2.71 \\
			200 & 5 & 2000 & 1405 & 0.14 & 1.43 & 0.01 & 1.79 & 0.01 & 0.25 & 0.01 & 0.17 & 0.01 & 0.11 & 0.01 & 7.84 \\
			\hline
			52 & 4  & 416 & 316 & 0.86 & 0.36 & 0.24 & 0.22 & 0.22 & 0.14 & 0.26 & 0.09 & 0.05 & 1.9 & - & - \\
			50 & 5  & 500 & 355 & 1.33 & 0.16 & 0.52 & 0.07 & 0.50 & 5.8 & 0.13 & 10.11 & 0.05 & 1.37 & 0.02 & 0.22 \\
			100 & 5 & 1000 & 705 & 0.41 & 0.33 & 0.19 & 8.58 & 0.04 & 6.14 & 0.02 & 0.41 & 0.02 & 0.58 & 0.02 & 3.78 \\
			200 & 5 & 2000 & 1405 & 0.14 & 1.55 & 0.07 & 39.07 & 0.04 & 9.44 & 0.01 & 0.12 & 0.02 & 4.48 & 0.01 & 6.66 \\
			\hline
		\end{tabular}
		\caption{Extended benchmark set, deterministic case, increasing number of costly servers}
		\label{table:extsetq}
	\end{table}
	
	The complexity of solving instances increases non-monotonically with the number of costly servers, for both increasing and decreasing inventory costs. While the 0-cost instances are comparatively easy to solve, the CPU time required to prove optimality increases with the addition of the first costly servers and then decreases as more servers are paid to store inventory. However, this is not true of the gap between the linear relaxation bound and the cost of the optimal solution, where the instances with no costly servers always have the largest gap.
	As expected, the CPU time required to prove optimality increases with the dimension, and an inverse correlation between storage cost and travel cost is more difficult to solve than the case with a positive correlation.
	
\end{APPENDICES}
	
\end{document}